\subsection{问题三的模型建立与求解}
\label{sec:5.4}

\subsubsection{模型建立}
\label{sec:5.4.1}

给定总预算 $C$（绝对 FLOPs）与外生上下文长度 $L_{ctx}$（token），在 $N$（个）、$D$（token）、$Q\in(0,1]$ 上最小化问题二的幂次主式。**主方案固定问题一在 A4 训练信赖域内选出的 17 域推荐配比 $\bp^*$**（赛题允许的建模选择，不是"配比与规模精确可分离"的定理），此时

\begin{equation*}
m_*=\exp\{s_{\rm red}\tilde\varphi(\bp^*)\}\approx0.8989423,\qquad
L=E+k_0(1-Q)+m_*\bigl(AN^{-\alpha}Q^{-\kappa_N}+BD^{-\beta}Q^{-\kappa_D}\bigr),
\end{equation*}

乘子只作用于两项随规模变化的损失，不作用于 $E$ 或地板；系数 $E,A,\alpha,B,\beta,\kappa_N,\kappa_D,k_0$ 全部读取 P2 接口，不在本问重拟合。约束为

\begin{equation*}
6ND+D\bigl[g(Q)-g(Q_0)\bigr]_++\eta NDL_{ctx}\le C,\qquad Q_0\le Q\le1,\qquad \eta=2\times10^{-4},
\end{equation*}

必须比较题目附录 B 的三类单位 token 提质成本：指数型 $g=10^7e^{6Q}$（凸）、幂函数型 $g=5\times10^9Q^4$（凸）、对数渐近型 $g=2\times10^9\ln(1+10Q)$（凹）。本问按\textbf{单次训练的静态预算情景}建模；题面未给训练起止日期、逐期预算、GPU 数、工期或碳排放等约束，不能把它们写成硬要求。题面建议 $C=10^{19},10^{22},10^{24}$ FLOPs，硬要求是至少三个不同量级；本项目自主采用 $10^{19}$ 至 $10^{25}$ 共七档网格，并把三个建议档单列。$L_{ctx}$ 是情景输入，不由优化器自由选择。连续 $N,D$ 与 $Q\ge Q_0$ 是为解析求解设定的建模假设。

\textbf{降维与闭式对照。} 三项成本都与 $D$ 成正比，而 $\partial L/\partial D=-m_*\beta T_D/D<0$，故在无额外 $D$ 上限的本模型内预算用尽。记 $K=6+\eta L_{ctx}$、$\Delta g=[g(Q)-g(Q_0)]_+$，由约束取等得

\begin{equation*}
D=\frac{C}{KN+\Delta g},
\end{equation*}

问题降为 $(\log_{10}N,Q)$ 的二维有界优化，仍须满足 $Q_0\le Q\le1$。固定 $Q$ 且令 $\Delta g=0$ 的解析对照为 $N^*=[\alpha AQ^{-\kappa_N}/(\beta BQ^{-\kappa_D})]^{1/(\alpha+\beta)}(C/K)^{\beta/(\alpha+\beta)}$、$D^*=C/(KN^*)$；式中 $m_*$ 与地板因对固定 $Q$ 的规模求导而抵消，其值仍影响不同 $Q$ 之间的最优选择。由此得

\begin{equation*}
\frac{D^*}{N^*}\propto\Bigl(\frac CK\Bigr)^{(\alpha-\beta)/(\alpha+\beta)}Q^{2(\kappa_N-\kappa_D)/(\alpha+\beta)},
\end{equation*}

两个指数在当前参数下约为 $0.09677$ 与 $0.42065$；这只是\textbf{无提质开销对照}，真实 $\Delta g>0$ 情景的比值与预算斜率须由数值最优解计算。上下文长度的解析临界值由注意力项与训练项之比 $\eta L_{ctx}/6$ 给出：$L_{ctx}^{\rm crit}=6/\eta=30000$ token，在固定 $L_{ctx}$ 下该比值与 $C,N,D$ 无关；C7 五档中 8192 在临界值以下、32768 在其以上，30000 本身不是 C7 的实测档位。

\textbf{结构性转移的定义与判据。} 本项目把"最优策略的质变"操作化为\textbf{全局最优解的质量约束活跃状态变化}：\texttt{LOWER}（$Q^*=Q_0$）、\texttt{INTERIOR}、\texttt{UPPER}（$Q^*=1$），并注明扫描预算区间、离散化精度与边界分类容差；若两个内点局部极小分支交换全局最优地位，即使仍同属 \texttt{INTERIOR} 也单列"全局分支切换"。对固定 $Q$ 的规模最优解，令 $S=KN+\Delta g$，剖面目标的质量导数与边际比为

\begin{equation*}
F_Q=-k_0-\frac{m_*}{Q}(\kappa_NT_N+\kappa_DT_D)+\frac{m_*\beta T_Dg'(Q)}{S},\qquad
\Phi(C;Q)=\frac{\Bigl[k_0+\frac{m_*}{Q}(\kappa_NT_N+\kappa_DT_D)\Bigr]S}{m_*\beta T_Dg'(Q)},
\end{equation*}

其中 $N,D$ 均在该\textbf{固定 $Q$} 下优化，不能代入其他质量点的全局规模解。内层解光滑且未触及搜索边界时，$\Phi(C;Q_0)\le1$、内点 $\Phi(C;Q)=1$、$\Phi(C;1)\ge1$ 分别是相应状态的\textbf{局部一阶必要条件}；$\Phi=1$ 的根不是全局状态切换的充分条件，须核对剖面曲线与端点并分别输出局部根与全局临界点，二分只用于已括根且已检查单次切换的区间。$Q_0\ge1-10^{-10}$ 时只有唯一质量状态，记 \texttt{NO\_QUALITY\_RANGE}，临界预算用 JSON \texttt{null}；搜索区间未括根返回 \texttt{NOT\_BRACKETED} 与 \texttt{null}，不返回上界。公式正确性已做单点有限差分核对：$C=10^{20}$、$L_{ctx}=2048$、指数型成本、$Q=0.8$ 处解析导数 $-0.1164515205405$ 与差分值 $-0.1164515204932$ 的相对误差约 $4.1\times10^{-10}$（对应 $\Phi=1.63892372$），见 \texttt{doc/Q3/新式子影响文档审查.md}。

\textbf{敏感性设计与输出要求。} 敏感性至少覆盖：三类 $g$、C7 五档 $L_{ctx}$、$\eta=10^{-4},2\times10^{-4},3\times10^{-4}$（题设固定 $2\times10^{-4}$，另两值只是本项目扰动情景）、P2 参数重抽样、主/锚点两种质量映射、P1 域质量与配比不确定性，以及固定配比乘子 \texttt{m=1}、H\_tot 与 P2 无地板参数对照。其中 H\_tot 在固定 $\bp^*$ 时对完整基线 Loss 乘一个与 $N,D,Q$ 无关的正常数，故其最优 $N,D,Q$ 与 \texttt{m=1} 对照相同、最优 Loss 按该常数缩放，须专门核验这一同配置关系。验收阈值：预算超限相对误差不大于 $10^{-9}$，三项份额和误差小于 $10^{-9}$，$N,D>0$ 且 $Q\in[Q_0,1]$，预算使用率与 1 之差小于 $10^{-9}$，无提质开销时闭式 $N^*$ 相对误差小于 $10^{-5}$，关键 18 组求解 Loss 不劣于独立粗网格且差不超过 $2\times10^{-5}$。\texttt{P3\_optimal\_config.csv} 的份额分母统一为\textbf{实际总消耗}，同时另报预算使用率，不得混用分母；\texttt{C,N,D} 均为绝对量。

\subsubsection{求解方法}
\label{sec:5.4.2}

问题三的求解器与接口脚本位于 \texttt{src/code\_q3/}（\texttt{q3\_00\_model.py}、\texttt{q3\_01\_optimize.py}、\texttt{q3\_02\_transition.py}、\texttt{q3\_03\_mixture\_joint.py}），联立配比探索性脚本的随机种子为 $2026$。**本版 Spec v5 的目标式（幂次 M4+地板并含 $m_*$）尚未在代码与 P3 接口中实现，因此 \ref{sec:5.4.3} 只给出求解方案与解析对照，数值结果待重跑；现有 \texttt{src/outputs\_q3/} 产物为旧线性口径，只作对照。

按 \ref{sec:5.4.1} 消去 $D$，在 $(\log_{10}N,Q)$ 上求解：内层对 $\log_{10}N$ 用黄金分割搜索（约束取等后的剖面损失关于 $N$ 单峰），外层对 $Q$ 用网格加密并保护端点 $Q_0$ 与 $1$，同时保留独立粗网格穷举核对，并分别输出局部 $\Phi=1$ 根与全局状态切换点。主网格为 3 种成本函数 $\times$ 7 档预算 $\times$ 5 档 C7 上下文长度 $\times$ 2 个同映射下的起始质量对照，共 210 组（本项目设计，不是赛题规定的组数）。临界预算的搜索用二分，且仅在已括根、已验证单次切换的区间内使用；未括根与 $Q_0=1$ 两种情形显式返回 \texttt{NOT\_BRACKETED}/\texttt{NO\_QUALITY\_RANGE} 与 \texttt{null}。

\subsubsection{结果与分析}
\label{sec:5.4.3}

\textbf{本节写作限制。} \texttt{doc/Q3/问题三\_SPEC.md} v5 已把目标式改为问题二的幂次 M4＋地板并保留固定配比乘子 $m_*\approx0.8989423$，但 \texttt{src/code\_q3/} 与 \texttt{src/outputs\_q3/} 尚未按该版重跑：现有产物生成于 2026-09-25 15:02，当时读取的是\textbf{线性质量式}的 $\kappa_N=0.4315$、$\kappa_D=0.1405$、$k_0=0.1407$ 与 $Q_0=0.6608889991$。因此本小节依次给出：规范目标下的解析结论；现有落盘结果的旧口径对照（明确不作结论）；待补证据清单。

\textbf{预算取等与降维。} 三项成本都与 $D$ 成正比，且 $\partial L/\partial D=-m_*\beta T_D/D<0$，故在无额外 $D$ 上限的本模型内预算用尽，$D=C/(KN+\Delta g)$，其中 $K=6+\eta L_{ctx}$、$\Delta g=[g(Q)-g(Q_0)]_+$；问题降为 $(\log_{10}N,Q)$ 的二维有界优化。

\textbf{上下文长度的解析临界值。} 注意力项与训练项之比为 $\eta L_{ctx}/6$，在固定 $L_{ctx}$ 下与 $C,N,D$ 无关；两者相等给出 $L_{ctx}^{\rm crit}=6/\eta=30000$ token。C7 的五档可行值中有 8192（临界值以下）与 32768（临界值以上）；30000 本身不是 C7 的实测档位，因此"成本交叉"与"预算驱动的最优策略转移"是两件事，须分开陈述。

\textbf{质量状态的分类与判据。} 质量约束的活跃状态分 \texttt{LOWER}（$Q^*=Q_0$）、\texttt{INTERIOR}、\texttt{UPPER}（$Q^*=1$）三类，离开下界与到达上界是两个不同的临界事件；$\Phi(C;Q_0)\le1$、内点 $\Phi(C;Q)=1$、$\Phi(C;1)\ge1$ 只是固定 $Q$ 意义下的\textbf{局部一阶必要条件}，全局切换可能发生在两个内点分支相等处。因此求解须同时输出局部根与全局临界点。该 $\Phi$ 表达式的正确性已有单点有限差分依据：$C=10^{20}$、$L_{ctx}=2048$、指数型成本、$Q=0.8$ 处解析导数 $-0.1164515205405$ 与差分 $-0.1164515204932$ 的相对误差约 $4.1\times10^{-10}$（对应 $\Phi=1.63892372$，见 \texttt{doc/Q3/新式子影响文档审查.md}）。这属于公式正确性证据，\textbf{不是任何配置或临界预算的结果}。

\textbf{配比乘子的处理。} 主情景按 P2 的 \texttt{p\_channel} 读取 $m_*=0.8989423$（由 \texttt{s\_red\_main=1.4996207320397785} 与 \texttt{phi\_tilde\_pstar=-0.07104223742472571} 给出），参考点是 $\bp_{\rm ref}=\bp_{\rm train\_mean}$ 而非 $\bp^*$；忽略配比效应的对照须单独标记 \texttt{m=1}。由于 P2 把跨规模配比通道标为 \texttt{exploratory}，含 $m_*$ 的结论\textbf{必须报告这一条件性来源}并给出 \texttt{m=1} 对照。固定 $\bp^*$ 时 H\_tot 对完整基线 Loss 乘一个与 $N,D,Q$ 无关的正常数，故其最优 $N,D,Q$ 应与 \texttt{m=1} 完全一致、最优 Loss 按该常数缩放——这一同配置关系是可核验的，而不是可优化出来的差异。

\textbf{现有落盘结果：旧口径对照（不作结论）。} \texttt{src/outputs\_q3/} 现有产物对应线性质量式与旧 $Q_0$。为便于后续核对，仅列出其结构而不作为本文结论。

\begin{table}[H]
\centering\small
\caption{现有 P3 产物的口径与状态计数（数据源：\texttt{interface/P3\_optimal\_config.csv}、\texttt{T1\_optimal\_all.csv}，）}
\label{tab:q3old}
\begin{tabularx}{\textwidth}{@{}p{3.7cm}X p{4.3cm}@{}}
\toprule
项目 & 现有落盘值 & 与 v5 目标规范的差异 \\
\midrule
质量式 & 线性 $1+\kappa(1-Q)$ & v5 要求幂次 $Q^{-\kappa}$ \\
$\kappa_N,\kappa_D,k_0$ & 0.4315, 0.1405, 0.1407 & 当前 P2 为 0.1416, 0.0112, 0.2280 \\
$Q_0$ & 0.6608889991 & 当前 P1/P2 为 0.6624116679870492 \\
配比乘子 $m_*$ & 未使用 & v5 要求 0.8989423 并另报 \texttt{m=1} \\
主网格状态计数 & 8 组 \texttt{Q0} 角点、18 组内点、184 组提满 & 须按 v5 重算 \\
临界预算（指数型/幂函数型/对数型） & 离开 $Q_0$ 为 18.874/19.399/18.166；到达 $Q=1$ 为 20.168/19.845/18.166 & 须按 v5 重算 \\
$\Phi=1$ 解析根 & 18.833/19.358/18.669 & 须按 v5 重算 \\
求解器 vs 粗网格 & 最大 Loss 差 $2.37\times10^{-5}$ & 已超出 v5 验收阈值 $2\times10^{-5}$ \\
闭式 $N^*$ 互验 & 最大相对误差 $1.89\times10^{-7}$ & 满足 $<10^{-5}$ \\
\bottomrule
\end{tabularx}
\end{table}

上表的三档预算列（$10^{19}/10^{22}/10^{24}$ FLOPs）与"七档网格（$10^{19}$–$10^{25}$）"都是本项目设计；题面只建议前三档并要求至少三个量级。现有产物还缺 P3 manifest、缺三项绝对成本与预算使用率字段、缺 \texttt{H\_tot}/\texttt{m=1}/映射情景与 P1 不确定性传播，且 \texttt{P3\_sensitivity\_Lctx.csv}、\texttt{P3\_structural\_transition.json} 的字段与 Spec v5 要求不一致（后者为 GBK 编码，下游读取有隐患）。\textbf{上述任何数值都不得被引用为本版结论。}

\textbf{待补证据清单（问题三）。}

1. 按 v5 目标式改写 \texttt{src/code\_q3/q3\_00\_model.py}（$h_N,h_D$、$\Phi$、闭式对照）后重跑全链，生成 \texttt{P3\_optimal\_config.csv}、\texttt{P3\_sensitivity\_Lctx.csv}、\texttt{P3\_structural\_transition.json} 与 \texttt{P3\_run\_manifest.json}（含 P1/P2/C7 输入 SHA-256）。\texttt{【需补证据：Q3-1、Q3-2、Q3-5】}
2. 主情景读取 $m_*$ 并另报 \texttt{m=1}、\texttt{H\_tot} 与 P2 无地板参数的对照；同一场景下比较四种口径的 $N,D,Q,L$。\texttt{【需补证据：Q3-6、Q3-9】}
3. 分别输出局部 $\Phi=1$ 根与全局状态切换点，并对 $Q_0=1$、未括根、多根三种情形做专用验收。\texttt{【需补证据：Q3-7】}
4. 补三项成本绝对值、总消耗与预算使用率字段，统一份额分母为实际总消耗，并把三个题面建议档单列。\texttt{【需补证据：Q3-8】}
5. P1 质量与配比误差、P2 锚点映射误差进入敏感性；缺少逐次重抽样接口时，把全链区间标注为未完成。\texttt{【需补证据：Q3-3、Q3-9】}
